% ============================================================
% Model: 3D 刚性转子与 Stark 效应
% ============================================================

\section{3D 刚性转子与 Stark 效应}
\label{sec:rigid-rotor-3d}

\begin{tipbox}[关键词标签]
刚性转子 · 球谐函数 · 转动光谱 · Stark效应 · 简并微扰论 · 选择定则 · 双原子分子
\end{tipbox}

% --------------------------------------------------------
% 1. 模型定义框
% --------------------------------------------------------
\begin{modelbox}[模型：三维空间中的刚性转子]
\textbf{物理情境}：双原子分子或线性多原子分子可近似为刚性转子——两原子核间距固定，整体绕质心转动。转动惯量 $I=\mu r_0^2$，$\mu$ 为约化质量，$r_0$ 为平衡核间距。

\textbf{哈密顿量}：
\[
\hat H_0=\frac{\hat{\mathbf{J}}^2}{2I},
\]
其中 $\hat{\mathbf{J}}$ 为总角动量算符，$[\hat J_i,\hat J_j]=i\hbar\varepsilon_{ijk}\hat J_k$。

\textbf{本征解}：
\begin{itemize}[nosep]
    \item 本征函数：球谐函数 $Y_{jm}(\theta,\varphi)$，$j=0,1,2,\ldots$，$m=-j,\ldots,j$
    \item 本征能量：$\displaystyle E_j^{(0)}=\frac{\hbar^2}{2I}j(j+1)$
    \item 简并度：$g_j=2j+1$（对 $m$ 的简并）
\end{itemize}

\textbf{物理参数}：
\begin{itemize}[nosep]
    \item 转动常数 $B=\hbar^2/(2I)$，常以波数单位表示 $\tilde B=B/(hc)$（单位 cm$^{-1}$）
    \item 典型值：HCl 的 $B\approx10.6$ cm$^{-1}$，对应转动光谱在远红外/微波区
\end{itemize}
\end{modelbox}

% --------------------------------------------------------
% 2. 关键公式框
% --------------------------------------------------------
\begin{formulabox}[核心公式]
\begin{enumerate}[label=\textbf{(\arabic*)}]
    \item \textbf{能级与跃迁}：
    \[
    E_j=Bj(j+1),\quad\Delta E=E_{j+1}-E_j=2B(j+1).
    \]
    \texteq{转动光谱线位置：$\tilde\nu=2\tilde B(j+1)$，等间距}

    \item \textbf{电偶极选择定则}：
    \[
    \Delta j=\pm1,\quad\Delta m=0,\pm1.
    \]
    \texteq{需要分子有永久电偶极矩（如 HCl、CO）；同核双原子分子无转动光谱}

    \item \textbf{Stark 微扰}（电场沿 $z$ 轴）：
    \[
    H'=-d\mathcal{E}\cos\theta,
    \]
    $d$ 为分子永久电偶极矩，$\mathcal{E}$ 为电场强度。

    \item \textbf{球谐函数递推}（关键矩阵元）：
    \[
    \cos\theta\,Y_{jm}=\sqrt{\frac{(j+1)^2-m^2}{(2j+1)(2j+3)}}\,Y_{j+1,m}
    +\sqrt{\frac{j^2-m^2}{(2j-1)(2j+1)}}\,Y_{j-1,m}.
    \]
    \texteq{选择定则 $\Delta j=\pm1$，$\Delta m=0$}

    \item \textbf{一级能量修正}：
    \[
    E_j^{(1)}=\langle jm|H'|jm\rangle=0\quad(\text{因 }Y_{jm}\text{ 有确定宇称，}\cos\theta\text{ 为奇}).
    \]

    \item \textbf{二级能量修正}：
    \[
    E_{jm}^{(2)}=\frac{Id^2\mathcal{E}^2}{\hbar^2}\cdot\frac{j(j+1)-3m^2}{j(j+1)(2j-1)(2j+3)}.
    \]
    \texteq{$j=0$ 时 $E_{00}^{(2)}=-Id^2\mathcal{E}^2/(3\hbar^2)$}
\end{enumerate}
\end{formulabox}

% --------------------------------------------------------
% 3. 训练点框
% --------------------------------------------------------
\begin{drillbox}[训练点与考法变体]
\trainingpoint{1} \textbf{简并微扰论的适用性分析}：$H'=-d\mathcal{E}\cos\theta$ 在简并子空间 $\{|jm\rangle\}_{m=-j}^j$ 内的矩阵元全为零（$\Delta j=0$ 时积分为零）。因此一级修正为零，二级修正可直接用非简并微扰公式。

\trainingpoint{2} \textbf{简并解除程度}：微扰后能量依赖于 $m^2$。相同 $j$、不同 $|m|$ 的态能量不同，但 $+m$ 与 $-m$ 仍简并（时间反演对称性）。$m=0$ 下移最多，$|m|=j$ 下移最少（或上移）。

\trainingpoint{3} \textbf{转动光谱计算}：给定分子的转动常数 $B$ 和电偶极矩 $d$，计算前几条转动光谱线的频率/波数，以及 Stark 分裂的大小。

\trainingpoint{4} \textbf{与平面转子的对比}：3D 转子一级 Stark 效应为零（二次效应），但某些特殊构型（如对称陀螺）可能出现一级效应。

\trainingpoint{5} \textbf{非刚性修正}：实际分子有离心畸变，$E_j=Bj(j+1)-Dj^2(j+1)^2$，$D\ll B$。高 $j$ 时能级略低于刚性转子预测。
\end{drillbox}

% --------------------------------------------------------
% 4. 解题技巧框
% --------------------------------------------------------
\begin{tipbox}[快速识别与解题技巧]
\begin{itemize}
    \item \examnote{识别标志}：题干出现``刚性转子''''双原子分子''''转动光谱''''球谐函数''''Stark效应''''微波''，立即定位本模型。
    \item \examnote{球谐函数速查}：
    \begin{center}
    \begin{tabular}{lll}
    \toprule
    $Y_{00}$ & $\sqrt{1/4\pi}$ & $s$ 波，球对称 \\
    $Y_{10}$ & $\sqrt{3/4\pi}\cos\theta$ & $p_z$ \\
    $Y_{1,\pm1}$ & $\mp\sqrt{3/8\pi}\sin\theta e^{\pm i\varphi}$ & $p_x\pm ip_y$ \\
    $Y_{20}$ & $\sqrt{5/16\pi}(3\cos^2\theta-1)$ & $d_{z^2}$ \\
    \bottomrule
    \end{tabular}
    \end{center}
    \item \examnote{能级公式记忆}：$E_j=Bj(j+1)$。$E_0=0$，$E_1=2B$，$E_2=6B$，$E_3=12B$。跃迁 $j\to j+1$：$\Delta E=2B(j+1)$，光谱线间距 $=2B$。
    \item \examnote{Stark 效应级数判定}：$H'\propto\cos\theta$ 是奇宇称算符，$Y_{jm}$ 有 $(-1)^j$ 宇称。$\langle jm|\cos\theta|jm\rangle=0$（相同宇称），故一级修正为零，二级修正 $\propto\mathcal{E}^2$。
    \item \examnote{与氢原子的类比}：3D 转子的 $H\propto J^2$ 与氢原子的 $L^2$ 结构相似，但转子是刚性的（无径向自由度），而氢原子有径向运动。
\end{itemize}
\end{tipbox}

% --------------------------------------------------------
% 5. 易错点框
% --------------------------------------------------------
\begin{errorbox}{常见失误与陷阱}
\begin{itemize}
    \item \textbf{转动常数的单位}：$B$ 常以波数单位 $\tilde B=B/(hc)$（cm$^{-1}$）给出。计算频率时需转换：$\nu=2c\tilde B(j+1)$。注意 $1\,\text{cm}^{-1}\approx29.98\,$GHz。
    \item \textbf{选择定则的适用范围}：$\Delta j=\pm1$ 仅对具有永久电偶极矩的分子成立。同核双原子分子（$\text{H}_2$、$\text{N}_2$、$\text{O}_2$）电偶极矩为零，无纯转动光谱。
    \item \textbf{简并微扰的适用范围}：$H'=-d\mathcal{E}\cos\theta$ 在简并子空间内矩阵元为零，但这不意味着简并被完全解除。二级修正后 $+m$ 和 $-m$ 仍简并（时间反演对称性）。
    \item \textbf{球谐函数的归一化}：$\int|Y_{jm}|^2d\Omega=1$。递推公式中的系数已包含归一化，直接使用即可，不要额外乘因子。
    \item \textbf{$j=0$ 的特殊性}：基态 $j=0$ 非简并，$m=0$。二级修正 $E_{00}^{(2)}=-Id^2\mathcal{E}^2/(3\hbar^2)$。注意 $j=0$ 的公式中分母有 $(2j-1)=-1$，需小心符号。
    \item \textbf{离心畸变}：高 $j$ 时分子键长略伸长，$I$ 略增大，$B$ 略减小。修正项 $-Dj^2(j+1)^2$ 使能级略低于刚性转子预测，光谱线间距略小于 $2B$。
\end{itemize}
\end{errorbox}

% --------------------------------------------------------
% 6. 物理图像与关联模型
% --------------------------------------------------------
\begin{insightbox}[物理图像与模型关联]
\textbf{物理图像}：

3D 刚性转子就像一个在空中自由旋转的哑铃：
\begin{itemize}[nosep]
    \item 它可以绕任意轴转动，角动量 $\mathbf{J}$ 是矢量。
    \item 角动量大小量子化：$|\mathbf{J}|=\sqrt{j(j+1)}\hbar$。
    \item 角动量在 $z$ 轴的投影也量子化：$J_z=m\hbar$，$m=-j,\ldots,j$。
    \item 转得越快（$j$ 越大），动能越高——$E\propto j(j+1)$。
    \item 加电场后，哑铃倾向于沿电场方向排列，但量子力学不允许它完全固定（不确定性原理），只能以概率分布偏向电场方向。
\end{itemize}

\textbf{与相似模型的对比}：
\begin{center}
\begin{tabular}{llll}
\toprule
\textbf{对比维度} & \textbf{3D 刚性转子} & \textbf{氢原子} & \textbf{谐振子} \\
\midrule
势场 & 无（约束在球面） & $-e^2/r$ & $\frac12m\omega^2r^2$ \\
能级 & $Ej(j+1)$ & $-1/n^2$ & $n\hbar\omega$ \\
简并来源 & $m$（转动对称） & $l$（库仑对称） & $n$（$SU(d)$ 对称） \\
波函数 & $Y_{jm}$ & $R_{nl}Y_{lm}$ & 厄米$\times$高斯 \\
Stark效应 & 二级 & 线性（$n\geq2$） & 线性 \\
\bottomrule
\end{tabular}
\end{center}

\textbf{推广方向}：
\begin{itemize}[nosep]
    \item 对称陀螺：三个主轴转动惯量不同，$H=A\hat J_a^2+B\hat J_b^2+C\hat J_c^2$，能级更复杂
    \item 非刚性转子：振动-转动耦合，$B_v=B_e-\alpha_e(v+1/2)$
    \item 微波激射器（maser）：氨分子 $(\text{NH}_3)$ 的反转跃迁，频率 23.87 GHz
\end{itemize}
\end{insightbox}

% --------------------------------------------------------
% 7. 推导附录
% --------------------------------------------------------
\begin{derivationbox}[推导细节（自我检验用）]
\textbf{Step 1：球谐函数递推关系}

利用升降算符或递推公式：
\[
\cos\theta\,Y_{jm}=\sqrt{\frac{(j+1)^2-m^2}{(2j+1)(2j+3)}}\,Y_{j+1,m}
+\sqrt{\frac{j^2-m^2}{(2j-1)(2j+1)}}\,Y_{j-1,m}.
\]

记系数为 $C_+$ 和 $C_-$。

\textbf{Step 2：二级能量修正}

\[
E_{jm}^{(2)}=\sum_{j'\neq j}\frac{|\langle j'm|H'|jm\rangle|^2}{E_j-E_{j'}}.
\]

非零矩阵元仅 $j'=j\pm1$：
\[
\langle j+1,m|H'|jm\rangle=-d\mathcal{E}\,C_+,\quad
\langle j-1,m|H'|jm\rangle=-d\mathcal{E}\,C_-.
\]

能量差：
\[
E_j-E_{j+1}=\frac{\hbar^2}{2I}\bigl[j(j+1)-(j+1)(j+2)\bigr]=-\frac{\hbar^2}{I}(j+1),
\]
\[
E_j-E_{j-1}=\frac{\hbar^2}{2I}\bigl[j(j+1)-(j-1)j\bigr]=\frac{\hbar^2}{I}j.
\]

代入：
\begin{align*}
E_{jm}^{(2)}&=d^2\mathcal{E}^2\Bigl(\frac{C_+^2}{E_j-E_{j+1}}+\frac{C_-^2}{E_j-E_{j-1}}\Bigr)\\
&=\frac{Id^2\mathcal{E}^2}{\hbar^2}\Bigl(-\frac{C_+^2}{j+1}+\frac{C_-^2}{j}\Bigr).
\end{align*}

将 $C_+$ 和 $C_-$ 代入并化简：
\[
E_{jm}^{(2)}=\frac{Id^2\mathcal{E}^2}{\hbar^2}\cdot\frac{j(j+1)-3m^2}{j(j+1)(2j-1)(2j+3)}.
\]

\textbf{Step 3：$j=0$ 验证}

$j=0$ 时 $m=0$，$C_+=\sqrt{1/3}$，$C_-=0$。

$E_{00}^{(2)}=\frac{Id^2\mathcal{E}^2}{\hbar^2}\cdot\frac{-1/3}{1}=-\frac{Id^2\mathcal{E}^2}{3\hbar^2}$。

与通式一致：$j(j+1)-3m^2=0$，分母 $(2j-1)(2j+3)=(-1)(3)=-3$，故
\[
E_{00}^{(2)}=\frac{Id^2\mathcal{E}^2}{\hbar^2}\cdot\frac{0}{0\cdot(-3)}\;(?)
\]

实际上 $j=0$ 时 $C_-$ 无定义（$j^2-m^2=0$ 但分母 $2j-1=-1$），需直接用两项公式，仅 $j\to j+1=1$ 项有贡献。
\end{derivationbox}

\vspace{1em}
